\SourceLeaf{271}%
\noindent folgt. Die Konstante ist jedoch gleich null zu setzen, da $y$ für $x = 0$ verschwindet, sodass schließlich

\[
y = \alpha \,.\, \frac{\sin kx}{k}
\]

gilt. Demnach ist

\[
f(x) = \alpha\, \frac{\sin kx}{k}
\]

und nach 14)

\[
\tag*{(G)} g(x) = \cos kx.
\]

17. Setzt man für $g(x)$ den Ausdruck $\cos kx$ in die Beziehungen V. 1), 2), 3) und 4) ein, so ergeben sich die 10 Beziehungen

\[
\cos kc = \cos ka \,.\, \cos kb,
\]

\[
\begin{array}{cc}
\mathrm{tg}\, kb = \mathrm{tg}\, kc \,.\, \cos A, & \mathrm{tg}\, ka = \mathrm{tg}\, kc \,.\, \cos B, \\[4pt]
\mathrm{tg}\, ka = \mathrm{tg}\, A \,.\, \sin kb, & \mathrm{tg}\, kb = \mathrm{tg}\, B \,.\, \sin ka, \\[4pt]
\sin ka = \sin A \,.\, \sin kc, & \sin kb = \sin B \,.\, \sin kc, \\[4pt]
\cos A = \cos ka \,.\, \sin B, & \cos B = \cos kb \,.\, \sin A,
\end{array}
\]

\[
\cos kc = \cot A \,.\, \cot B,
\]

und das sind die bekannten Gleichungen der \textls{sphärischen Trigonometrie}, wenn man den Radius der Kugel gleich $\dfrac{1}{k}$ setzt. Sie gehen in die Gleichungen der \textls{nichteuklidischen Trigonometrie} über, wenn man für $k$ einen rein imaginären Wert $k = ik'$ wählt.

Es fällt auf, dass \textsc{Gauss} die Konstante, die bei der Integration der Differentialgleichung 15) auftritt, mit $k$ bezeichnet. Ob er damit einen Zusammenhang mit dem Krümmungsmaß andeuten wollte, das er in den Disquisitiones generales circa superficies curvas ebenfalls $k$ nennt, lässt sich freilich nicht mit Bestimmtheit behaupten. In diesem Zusammenhang verdient eine Äußerung von \textsc{Minding} in seiner Abhandlung \emph{Beiträge zur Theorie der kürzesten Linien auf krummen Flächen} Erwähnung, die 1840 in \textsc{Crelles} Journal für die reine und angewandte Mathematik (Bd. 20, S. 323--327) erschienen ist:

»Dass auf jeder Fläche von unveränderlichem positivem Krümmungsmasse zwischen den Seiten und Winkeln eines aus kürzesten Linien gebildeten Dreiecks die Formeln der sphärischen Trigonometrie gelten, folgt sogleich, wenn man sich erinnert, dass jede Fläche dieser Art sich auf eine Kugel abwickeln lässt. Ist das Krümmungsmass negativ, so gelten dieselben Formeln mit der Änderung, dass die hyperbolischen Functionen der Seiten an die Stelle der trigonometrischen treten. Sind nemlich $a, b, c$ die Seiten des Dreiecks, $A$ der Gegenwinkel von $a$ und $k$ das unveränderliche Krümmungsmass, gleichviel ob positiv oder negativ, so ist es nicht schwer, die Richtigkeit folgender Gleichung zu beweisen:

\[
\cos a\BareRoot{k} = \cos b\BareRoot{k} \,.\, \cos c\BareRoot{k} + \sin b\BareRoot{k} \,.\, \sin c\BareRoot{k} \,.\, \cos A.\text{«}
\]

Die Rotationsflächen von konstantem negativem Krümmungsmaß hatte \textsc{Minding} bereits 1839 bestimmt (\textsc{Crelles} Journal, Bd. 19). Dazu vergleiche man die folgende Notiz von \textsc{Gauss}, die in einem Handbuch steht und aus der Zeit zwischen 1823 und 1827 stammt:
